%%
%% This is file `math-operator_heading.tex',
%% generated with the docstrip utility.
%%
%% The original source files were:
%%
%% math-operator_code.dtx  (with options: `heading')
%% 
%% This file is from version 1.0 of the free and open-source
%% LaTeX package "math-operator," released February 2025.
%% 
%% Copyright 2025 Conrad Kosowsky
%% 
%% This file may be used, distributed, and modified under the
%% terms of the LaTeX Public Project License, version 1.3c or
%% any later version. The most recent version of this license
%% is available online at
%% 
%%           https://www.latex-project.org/lppl/
%% 
%% This work has the LPPL status "maintained," and the current
%% maintainer is the package author, Conrad Kosowsky. He can
%% be reached at kosowsky.latex@gmail.com.
%% 
%% PLEASE KNOW THAT THIS FREE SOFTWARE IS PROVIDED WITHOUT
%% ANY WARRANTY. SPECIFICALLY, THE "NO WARRANTY" SECTION OF
%% THE LATEX PROJECT PUBLIC LICENSE STATES THE FOLLOWING:
%% 
%% THERE IS NO WARRANTY FOR THE WORK. EXCEPT WHEN OTHERWISE
%% STATED IN WRITING, THE COPYRIGHT HOLDER PROVIDES THE WORK
%% `AS IS’, WITHOUT WARRANTY OF ANY KIND, EITHER EXPRESSED
%% OR IMPLIED, INCLUDING, BUT NOT LIMITED TO, THE IMPLIED
%% WARRANTIES OF MERCHANTABILITY AND FITNESS FOR A PARTICULAR
%% PURPOSE. THE ENTIRE RISK AS TO THE QUALITY AND PERFORMANCE
%% OF THE WORK IS WITH YOU. SHOULD THE WORK PROVE DEFECTIVE,
%% YOU ASSUME THE COST OF ALL NECESSARY SERVICING, REPAIR, OR
%% CORRECTION.
%% 
%% IN NO EVENT UNLESS REQUIRED BY APPLICABLE LAW OR AGREED
%% TO IN WRITING WILL THE COPYRIGHT HOLDER, OR ANY AUTHOR
%% NAMED IN THE COMPONENTS OF THE WORK, OR ANY OTHER PARTY
%% WHO MAY DISTRIBUTE AND/OR MODIFY THE WORK AS PERMITTED
%% ABOVE, BE LIABLE TO YOU FOR DAMAGES, INCLUDING ANY GENERAL,
%% SPECIAL, INCIDENTAL OR CONSEQUENTIAL DAMAGES ARISING OUT
%% OF ANY USE OF THE WORK OR OUT OF INABILITY TO USE THE WORK
%% (INCLUDING, BUT NOT LIMITED TO, LOSS OF DATA, DATA BEING
%% RENDERED INACCURATE, OR LOSSES SUSTAINED BY ANYONE AS A
%% RESULT OF ANY FAILURE OF THE WORK TO OPERATE WITH ANY
%% OTHER PROGRAMS), EVEN IF THE COPYRIGHT HOLDER OR SAID
%% AUTHOR OR SAID OTHER PARTY HAS BEEN ADVISED OF THE
%% POSSIBILITY OF SUCH DAMAGES.
%% 
%% For more information, see the LaTeX Project Public License.
%% Derivative works based on this software may come with their
%% own license or terms of use, and the package author is not
%% responsible for any third-party software.
%% 
%% Happy TeXing!
%% 

\csname count@\endcsname\catcode`\@
\makeatletter

\def\packagedate{February 2025}
\def\packageversion{1.0}

\let\@@section\section
\def\section{\relax
  \vskip1in\kern\z@\vskip-1in\vskip\z@
  \vskip1.5in\kern\z@\vskip-1.5in\vskip\z@
  \@ifstar\star@sect\no@star@sect}
\def\star@sect#1{\@@section*{#1}\section@name{#1}}
\def\no@star@sect#1{\@@section{#1}\section@name{#1}}
\def\section@name#1{\expandafter
  \def\csname section@\thesection\endcsname{#1}}
\def\sectionname{\csname section@\thesection\endcsname}

\protected\def\clap#1{\hb@xt@\z@{\hss#1\hss}}
\def\@oddhead{\ifnum\count0>1\relax
  \rlap{}\hfil\clap{\itshape\sectionname}\hfil
  \llap{\the\count0}\fi}
\def\@evenhead{\ifnum\count0>1\relax
  \rlap{\the\count0}\hfil\clap{\itshape\sectionname}\hfil
  \llap{}\fi}
\def\@oddfoot{\hfil\ifnum\count0=1\relax1\fi\hfil}
\let\@evenfoot\@empty

\pretolerance=20
\hyphenpenalty=10
\exhyphenpenalty=5
\brokenpenalty=0
\clubpenalty=5
\widowpenalty=5
\finalhyphendemerits=300
\doublehyphendemerits=500

\flushbottom

\def\topfraction{1}
\def\bottomfraction{1}
\let\code\@undefined
\newenvironment{code}
  {\strut\vadjust\bgroup\vskip 5pt plus 1pt minus 3pt\relax
    \parindent\z@\leftskip2em\relax
    \noindent\strut\ignorespaces}
  {\strut\par\vskip 5pt plus 1pt minus 3pt\relax
    \egroup\hfill\break\strut\ignorespacesafterend}
\def\vrb#1{\expandafter\texttt\expandafter{\string#1}}
\parskip=0pt

{\large
  \parindent=0pt\relax
  \leftskip=0pt plus 1fill\relax
  \rightskip=0pt plus 1fill\relax
{\strut\Large Package \textsf{math-operator} v.\ \packageversion\ \documentname\par
{\strut Conrad Kosowsky}\par
{\strut\packagedate}\par
{\strut\ttfamily kosowsky.latex@gmail.com}\par}

\bigskip

{\small
\leftskip=0.5in\relax
\rightskip=0.5in\relax

\centerline{\bfseries Overview\strut}
\noindent The \textsf{math-operator} package defines control sequences for roughly one hundred and fifty math operators, including special functions, probability distributions, pure mathematical constructions, and a variant of |\overline|. The package also provides an interface for users to define new math operators similar to the \textsf{amsopn} package. New operators can be medium or bold weight, and they may be declared as |\mathord| or |\mathop| subformulas.\par}

\bigskip\bigskip\nointerlineskip
\centerline{\vrule height 0.5pt width 2.5in}\bigskip\bigskip}
\nointerlineskip

\catcode`\@\count@

\endinput
%%
%% End of file `math-operator_heading.tex'.
